\stopchapter

\chapter{Die Macht der Sprache}

\section{2024-11-13}

\startitemize[packed]
\item wiederholung
\item die Sprache Lehrt mich neu zu kennen
\item viele ähnliche wörter
\item 
\stopitemize

\section{2024-11-19}

Wolf Schneider

\startitemize[packed]
\item Wortfälschung -> Scheinbarjetzt
\item Anglizismen
\item Internet vermehrt gesprochene Worte (mehr Kommunikation), aber abnehmende Qualität => Sorgfalt
\item "Kietzdeutsch" / "Sprache mit migrationshintergrund"
\stopitemize

\subsection{Kritik am Sprachwandel}

\startitemize[packed]
\item Sprache als Waffe: SPrache kann ausgrenzen, verletzen oder töten
\item Freund der Sprache: vermeintlich gebildiete Schichten, Traditionalisten der deutschen Sprache (Sprachbewahrere)
\stopitemize

\subsection{Standardisierungsideologie vs Differenztheorie}

\startplacetable[location={here},title={Vergleich}]
\setupTABLE[r][1][background=color,backgroundcolor=gray]
\bTABLE[frame=off]
\startCSV
Standardisierungsideologie, Differenztheorie
\stopCSV
\bTR
\bTD % Standardisierungsideologie
\startitemize[packed]
\item schriftliche Standardsprache (DUDEN) sollte der maßstab sein
\item "schlechtes Deutsch" alles was davon abweicht
\stopitemize
\eTD
\bTD % Differenztheorie
\startitemize[packed]
\item Sprache wird nicht als "gut" oder "schlecht" beurteilt
\item alle Varietäten sind funktional grundsätzlich gleichwertig
\stopitemize
\eTD
\eTR
\eTABLE
\stopplacetable

\subsection{Sprache: Politisch-Geselschaftlich}

\startitemize[packed]
\item Hitlerreden
\item "Vorwärts immer, rückwärts nimmer"
\item Lang + Ausgeschmückt
\item viele Worte, wenig Inhalt
\item Populismus
\stopitemize
